\subsection{Desenvolvimento das operações de acoplamento}
\label{subsec_desenvolvimento_acoplamento}

Após as primeiras demonstrações de \textit{rendezvous} orbital, o desenvolvimento das operações de acoplamento passou a concentrar esforços na criação de sistemas capazes de realizar a aproximação final e o contato mecânico entre veículos espaciais de forma segura e repetível. 
Esse avanço envolveu tanto a evolução de mecanismos de acoplamento quanto o desenvolvimento de sistemas de navegação relativa, guiamento e controle capazes de operar em distâncias progressivamente menores entre as espaçonaves.

Nas primeiras aplicações, as operações de acoplamento dependiam fortemente da atuação de pilotos humanos, responsáveis por interpretar informações de navegação relativa e conduzir manualmente as fases finais da aproximação. 
Com o avanço da tecnologia de sensores e de sistemas embarcados, passaram a ser desenvolvidas arquiteturas capazes de automatizar parte dessas operações, reduzindo a carga operacional sobre os astronautas e aumentando a previsibilidade do processo de aproximação.

Um dos desafios associados a essa evolução consiste em garantir a estimativa precisa da posição e orientação relativas entre os veículos durante a fase final de aproximação. 
Nesse contexto, diferentes soluções foram propostas ao longo do tempo, incluindo sistemas baseados em sensores ópticos, sensores a laser e técnicas de fusão de dados para estimar o estado relativo entre perseguidor e alvo.

Estudos iniciais sobre acoplamento automático entre espaçonaves investigaram arquiteturas de sensoriamento e controle capazes de conduzir o processo sem intervenção humana direta. 
Um exemplo é o trabalho apresentado por \cite{micheal1981}, no qual foi analisado um cenário de acoplamento associado ao projeto \emph{Teleoperator Retrieval System (TRS)}. 
Nesse estudo, foram avaliados métodos para estimativa de posição e atitude relativas utilizando sensores ópticos e sistemas de varredura a laser, evidenciando a importância da integração entre sensoriamento e algoritmos de controle para viabilizar operações autônomas de acoplamento.

Com a evolução das missões de logística orbital e de suporte a estações espaciais, sistemas de acoplamento altamente automatizados passaram a ser empregados em veículos de serviço. 
Um exemplo representativo é o \emph{Automated Transfer Vehicle} (ATV), desenvolvido pela \gls{esa} para missões de reabastecimento da \gls{iss}. 
O ATV realizava operações de \gls{rvdb} utilizando sistemas automáticos de navegação relativa, integrando sensores, câmeras e sistemas de guiamento para ajustar continuamente sua trajetória durante a fase final de aproximação \cite{esaATVRendezvous}.

Durante a aproximação final, o veículo utilizava sensores a laser para estimar a distância e o alinhamento em relação à estação espacial, permitindo realizar correções de trajetória em tempo real. 
A automação dessas operações reduziu significativamente a necessidade de intervenção humana direta e demonstrou a viabilidade de sistemas de \gls{rvdb} altamente automatizados em missões operacionais.